\section{Finite descent, modular gluing, and the boundary of inverse limits}
\label{sec:descent}

The categorical layer becomes useful when it identifies the exact local
information required to assemble a mathematical object. This section gives
two positive gluing results, an explicit obstruction, and a constructive
resolution under a stronger local lifting hypothesis. All inputs are numerical
tables or fixed PR families of such tables.

\subsection{A finite sheaf of function tables}

Let $X=[n]$, let $U_0,\ldots,U_{k-1}$ be subsets represented by Boolean masks,
and let the codomain be $[m]$. A section on $U_i$ is a finite table
$f_i:U_i\to[m]$. Suppose the masks cover $X$ and the sections agree on
overlaps:
\begin{align}
 &\forall x<n\ \exists i<k\ (x\in U_i),\label{eq:finite-cover}\\
 &x\in U_i\cap U_j\quad\Longrightarrow\quad f_i(x)=f_j(x).
 \label{eq:overlap}
\end{align}
Both assertions have bounded Boolean tests. A finite table $f:X\to[m]$
restricts to compatible sections on every such cover.

\begin{theorem}[Uniform finite gluing]
\label{thm:finite-glue}
There is a PR algorithm which, from data satisfying
\eqref{eq:finite-cover}--\eqref{eq:overlap}, constructs the unique table
$f:X\to[m]$ restricting to every $f_i$. The construction and its
uniqueness are provable in iPRA. The correspondence between global sections
and compatible local sections is a PR bijection.
\end{theorem}

\begin{proof}
For $x<n$, let $i(x)$ be the first mask containing $x$, obtained by a scan of
the $k$ masks. Coverage proves that the scan succeeds. Define
$f(x)=f_{i(x)}(x)$. If $x\in U_j$, overlap compatibility gives
$f_{i(x)}(x)=f_j(x)$, so $f$ has the required restrictions.
Any other global table $g$ with those restrictions satisfies
$g(x)=f_{i(x)}(x)=f(x)$ at each $x<n$, proving uniqueness by finite
extensionality. Restriction followed by this gluing recovers a global table
entry by entry. Gluing followed by restriction recovers each local table by
the overlap equation. The two maps are therefore inverse.

The selector uses a bounded first-success register and the output uses a
loop of length $n$. The invariant says that every output entry constructed
so far agrees with every supplied section defined there. It is a bounded
conjunction over at most $nk$ pairs; Lemma~\ref{lem:booleanization} turns it
into a QF invariant. The empty case $n=0$ yields the empty table. If $m=0$
and $n>0$, the hypotheses cannot hold, since some nonempty covered domain
would have to map to the empty set.
\end{proof}

In categorical terms, restriction is an equalizer presentation:
\[
 \operatorname{Map}(X,[m])\ \cong\
 \operatorname{Eq}\!\left(
  \prod_i\operatorname{Map}(U_i,[m])
  \rightrightarrows
  \prod_{i,j}\operatorname{Map}(U_i\cap U_j,[m])\right).
\]
The two arrows restrict the first and second member of an overlapping pair.
All products and function sets here are finite and were constructed explicitly
in the preceding section. This is a finite sheaf condition, with a proved
gluing operation rather than an axiom asserting that a global section exists.

It also respects changes of cover. Given a refinement $V_j\subseteq
U_{a(j)}$, restrict the old sections to the new masks. If the refined masks
still cover $X$, both old and new gluing give the same table, because their
restrictions agree at every point. Composing an output table with a fixed
map $[m]\to[q]$ also commutes with gluing, by evaluating at $i(x)$.
These two naturality statements are finite pointwise equalities. Thus a
choice of the first covering patch affects the implementation but not the
represented result.

\begin{example}
Take $X=[5]$, $U_0=\{0,1,2\}$, $U_1=\{2,3,4\}$, and values in $[7]$.
The local tables $(0,1,4)$ and $(4,2,2)$ agree at the overlap $2$.
Their unique gluing is $(0,1,4,2,2)$, the table of $x^2\bmod7$.
If one overlap value is changed, the algorithm detects a failed compatibility
equation. It does not manufacture a section from inconsistent local data.
\end{example}

For finite observations of an analytic object, exact overlap equality may be
unavailable. One may replace it by a bounded relation of compatible error
intervals and compute a finite matching family, using the existing
finite-demand methods \cite{repo-evidence,repo-demand}. The theorem above
should not be applied by replacing exact equality with a tolerance and
retaining the same uniqueness conclusion. Approximate compatibility is a
different contract whose error propagation must be proved.

\subsection{A modular pullback with an explicit inverse}

For positive integers $a,b$, put $g=\gcd(a,b)$ and $\ell=ab/g$.
Represent $\Z/a\Z$ by its canonical residues $[a]$. Reduction modulo $g$
gives a cospan, and its finite pullback is
\[
 P=\{(r,s):r<a,\ s<b,\ r\equiv s\pmod g\}.
\]
The map $\theta:[\ell]\to P$ sends $x$ to $(x\bmod a,x\bmod b)$.
This example shows how a local compatibility condition gives an actual
global constructor.

\begin{theorem}[Constructive two-modulus Chinese remainder isomorphism]
\label{thm:crt}
The map $\theta$ is a PR bijection uniformly in $a,b>0$, with PR inverse
provable correct in iPRA. Pairwise coprimality is unnecessary; agreement
modulo $g$ is exactly the compatibility condition.
\end{theorem}

\begin{proof}
Bounded extended Euclid produces integers $u,v$ such that
$ua+vb=g$. Its remainder sequence decreases whenever it is nonzero, so at
most $a+b+1$ updates suffice; the two tracked linear-combination identities
are bounded arithmetic invariants. Write $a'=a/g$, $b'=b/g$.
Then $ua'+vb'=1$.

For $(r,s)\in P$, let $q=(s-r)/g\in\Z$ and choose the canonical residue
$t=(uq)\bmod b'$. Define
\[
 x=(r+at)\bmod\ell.
\]
Since $a'u\equiv1\pmod {b'}$, one has
$a't\equiv q\pmod {b'}$, hence $at\equiv s-r\pmod b$.
Thus $x\equiv r\pmod a$ and $x\equiv s\pmod b$.
This proves surjectivity with the displayed constructor.

If $x,y$ have the same residues, let $z=x-y$. Both $a$ and $b$ divide $z$.
Write $z=ak$. From $b\mid ak$ follows $b'\mid a'k$.
Multiplying $ua'+vb'=1$ by $k$ shows that $b'\mid k$.
Consequently $\ell=ab'$ divides $z$. If $0\le x,y<\ell$, their
difference has absolute value below $\ell$; divisibility then forces $z=0$.
The constructed $x$ is therefore the unique representative, and the two
maps are inverse. All arithmetic, gcd, signed division by a known divisor,
and residue operations are fixed PR operations; the preceding argument uses
only their finite arithmetic contracts.
\end{proof}

For example, residues $5\bmod12$ and $11\bmod18$ agree modulo $6$.
The inverse gives $29\bmod36$. In diagrammatic form the theorem says that
the square
\[
\begin{tikzcd}[column sep=large,row sep=large]
 \Z/\ell\Z \arrow[r] \arrow[d] & \Z/a\Z \arrow[d] \\
 \Z/b\Z \arrow[r] & \Z/g\Z
\end{tikzcd}
\]
is a pullback in the finite table category. The result assembles compatible
finite residues. It says nothing about whether an infinite compatible
congruence tower is the expansion of a rational or ordinary integer solution
of a separate equation. That remains a distinct adequacy problem.

\subsection{Every finite horizon may be solvable without a computable branch}

The following direct construction makes the boundary precise. Fix a standard
enumeration $\varphi_e$ of partial computable functions and a PR bounded
Turing-machine simulation predicate. Concretely, finite configuration codes
have fixed PR initialization, one-step, halting, and output functions, and
\[
 \operatorname{run}(e,x,0)=\operatorname{init}(e,x),\qquad
 \operatorname{run}(e,x,t+1)
   =\operatorname{step}(e,\operatorname{run}(e,x,t)).
\]
Halted configurations are absorbing. Determinism and persistence of a
halted output are finite equations proved by QF induction on the run.
A binary word of length $n$ is a number
$c<2^n$, with bit $\sigma(e)$ at position $e<n$. Declare it admissible when
\[
 \begin{array}{ll}
 \sigma(e)=1,&\text{if }\varphi_e(e)\text{ has halted with output }0
                      \text{ within }n\text{ steps},\\
 \sigma(e)=0,&\text{if }\varphi_e(e)\text{ has halted with output }1
                      \text{ within }n\text{ steps}.
 \end{array}                                             \tag{T}
\]
Other outputs impose no condition. The simulations are deterministic, so
the two constraints never conflict. The finite admissibility test is PR.

\begin{theorem}[Finite solvability does not give computable coherence]
\label{thm:tree}
There is a uniformly PR sequence of nonempty finite sets $A_n$ of binary
words of length $n$, with PR restriction maps $A_{n+1}\to A_n$ and a PR
selector choosing one element of each $A_n$, but no computable binary
sequence whose length-$n$ prefix belongs to $A_n$ for every $n$.
\end{theorem}

\begin{proof}
Let $A_n$ consist of the words satisfying (T). Choose bit one when output
zero has been seen by time $n$, and bit zero otherwise. This produces a
member of $A_n$ by a loop of $n$ bounded simulations. Its uniform
correctness is a bounded, hence QF-inductible, assertion. Thus even the
levelwise witness producer is PR, not merely an unproved existence claim.

If a word of length $n+1$ satisfies (T), its length-$n$ prefix does also:
every computation seen within $n$ steps is still seen within $n+1$ steps
with the same output. Prefix truncation is the required PR restriction.

Suppose a total computable binary sequence $h$ had every prefix admissible.
Take an index $d$ computing $h$, and let $t$ be the halting time of
$\varphi_d(d)=h(d)$. For $n>\max(d,t)$, the admissibility of the
length-$n$ prefix forces $h(d)=1$ if its computed value is zero, and
$h(d)=0$ if its computed value is one. Both possibilities are contradictory.
No such computable coherent sequence exists.
\end{proof}

The nonexistence conclusion is a metatheoretic computability statement;
the construction and levelwise witnesses are explicit arithmetic data.
It does not say that the finite limits constructed earlier fail, that no
classical infinite branch exists, or that the bonding maps are surjective.
It says that one cannot replace all finite observation problems by a
computably available inverse-limit point without additional information.
The bounded interpreter used here is consistent with the impossibility
of an \emph{unbounded total} PR evaluator for all PR programs.

\subsection{The missing positive resource: finite liftings}

There is a useful constructive repair. Suppose a fixed PR family gives
finite sets
\[
 A_n=\{a<b(n):V(n,a)=1\}
\]
and bonding maps $p(n,-):A_{n+1}\to A_n$, with iPRA proofs of their
validity. The bounds $b$ and tests $V$ are fixed PR functions, not
uncoded families. Suppose a member $a_0\in A_0$ is supplied, and iPRA
proves
\begin{equation}
 a<b(n)\wedge V(n,a)=1\ \longrightarrow\
 \exists c<b(n+1)\bigl(V(n+1,c)=1\wedge p(n,c)=a\bigr).
 \label{eq:surjective-lifting}
\end{equation}

\begin{theorem}[Coherent points from bounded surjective liftings]
\label{thm:tower}
Under these hypotheses a fixed PR function $a(n;a_0)$, with the supplied
$a_0$ retained as a numerical parameter, produces a coherent
point: writing $a(n)$ for $a(n;a_0)$, one has $a(n)\in A_n$ and
$p(n,a(n+1))=a(n)$ for every $n$.
These assertions have iPRA derivations from the displayed contracts.
\end{theorem}

\begin{proof}
Let $s(n,a)$ be the first $c<b(n+1)$ satisfying the conjunction in
\eqref{eq:surjective-lifting}, with a default if no such $c$ exists.
This is a bounded PR scan. The first-success lemma and the surjectivity
proof imply that, for valid $a$, it returns a valid lift and
$p(n,s(n,a))=a$. Define by ordinary primitive recursion
\[
 a(0)=a_0,\qquad a(n+1)=s(n,a(n)).
\]
Induct on the Boolean validity equation for $a(n)$, with its numerical
bound included as a PR comparison. The initial assertion is the supplied
one, and the local lifting contract proves the successor assertion.
Instantiate that contract once more at the resulting $a(n)$ to obtain
the adjacent coherence equation. Generalizing $n$ completes the proof.
There is no induction on an unbounded assertion that a whole tail exists.
\end{proof}

This explains exactly how some coherent analytic or algebraic towers are
available in iPRA. A valid local lifting algorithm, or bounded surjectivity
of the finite stages, is enough. Merely knowing that each finite stage is
inhabited is weaker. The counterexample in Theorem~\ref{thm:tree} lacks
the surjective lifting contract. Existing filtered-lifting constructions
\cite{repo-coherent} provide stronger local data in their particular rings;
they are examples of the positive mechanism, not applications of an
unrestricted inverse-limit principle.
